\lyseis{4.3}{Παραδείγματα επίλυσης προβλημάτων}
\sol{1} $3x+3=111$, $3x=108$, $x=36$: οι αριθμοί $36,\,37,\,38$.
\sol{2} $x+2x=51$, $3x=51$, $x=17$: οι αριθμοί $17$ και $34$.
\sol{3} $x+(x+2)+(x+4)+(x+6)=100$, $4x+12=100$, $x=22$: οι $22,\,24,\,26,\,28$.
\sol{4} $\frac{3}{5}x=36$, $x=36:\frac{3}{5}=36\cdot\frac{5}{3}=60$.
\sol{5} $x-\frac{x}{3}-4=10$, $\frac{2x}{3}=14$, $x=21$ €. Έλεγχος: $21-7-4=10$.
\sol{6} $x+(x+5)=29$, $2x=24$, $x=12$ (η Πολυξένη) και ο αδελφός $17$.
\sol{7} $250\cdot8+100x=4\,800$, $100x=2\,800$, $x=28$ ημέρες.
\sol{8} $\frac{6\cdot5}{2}=15$ αγώνες. Με δέκα ομάδες $\frac{10\cdot9}{2}=45$.
\sol{9} $\frac{3}{8}x=1{,}2$, $x=1{,}2\cdot\frac{8}{3}=3{,}2$ λίτρα.
\sol{10} Η αμοιβή είναι $6x$ και το μισό της $270$: $3x=270$, $x=90$ € ανά ημέρα.
\sol{11} $3{,}5+0{,}8x=13{,}1$, $0{,}8x=9{,}6$, $x=12$ χλμ.
\sol{12} Πλησιάζουν με $8+12=20$ χλμ. την ώρα, άρα $x=5:20=0{,}25$ ώρες $=15$ λεπτά.
\sol{13} Μέτρησε κάθε χειραψία δύο φορές. Σωστά: $\frac{12\cdot11}{2}=66$.
\sol{14} $x+\frac{x}{2}=36$, $\frac{3x}{2}=36$, $x=24$. Έλεγχος: $24+12=36$.
\sol{15} Τα ζευγάρια $(0,10),\,(1,9),\dots,(10,0)$: $11$ ζευγάρια. Έχει πολλές λύσεις γιατί έχουμε δύο αγνώστους και μία μόνο ισότητα.
\sol{16} Κάθε ομάδα παίζει δύο φορές με καθεμία από τις άλλες $7$: ο συνολικός αριθμός είναι $8\cdot7=56$ αγώνες (χωρίς διαίρεση με το $2$, γιατί ο αγώνας εντός και ο εκτός είναι διαφορετικοί).
\sol{17} Δέκα ζευγάρια με άθροισμα $21$ το καθένα ($1+20$, $2+19$, \dots). Άθροισμα $10\cdot21=210$.
\sol{18} Ελεύθερη απάντηση. Π.χ. οι αριθμοί τριών γειτονικών σπιτιών έχουν άθροισμα $93$: $x+(x+1)+(x+2)=93$, $x=30$.
\sol{19} Ελεύθερη απάντηση. Ο Διόφαντος ονόμαζε τον άγνωστο «αριθμό».
